 \documentclass[11pt,a4paper]{article}

%====================== PACKAGES ======================

\usepackage[french]{babel}
\usepackage{charter}
\usepackage[utf8x]{inputenc}
\usepackage{lastpage}
%pour gérer les positionnement d'images
\usepackage{float}
\usepackage{amsmath}
\usepackage{graphicx}
\usepackage[table]{xcolor}
\usepackage[colorinlistoftodos]{todonotes}
\usepackage{url}
%pour les informations sur un document compilé en PDF et les liens externes / internes
\usepackage{hyperref}
%pour la mise en page des tableaux
\usepackage{array}
\usepackage{tabularx}
%pour utiliser \floatbarrier
%\usepackage{placeins}
%\usepackage{floatrow}
%espacement entre les lignes
\usepackage{setspace}
%modifier la mise en page de l'abstract
\usepackage{abstract}
%police et mise en page (marges) du document
\usepackage[T1]{fontenc}
\usepackage[top=2cm, bottom=2cm, left=2cm, right=2cm]{geometry}
%Pour les galerie d'images
\usepackage{subfig}
%Pour les tirets stylisé (itemize)
\usepackage{enumitem}
\usepackage{pifont}
\usepackage{amsmath}
\usepackage{amssymb}
\usepackage{stmaryrd}


\usepackage{algorithm}
\usepackage{algpseudocode}
  
%Includes "References" in the table of contents
\usepackage[nottoc]{tocbibind}
\usepackage[maxnames=15, sorting=nty]{biblatex}
\addbibresource{bibliographie.bib}

%Pour le code
\usepackage{tcolorbox,listings}
\usepackage{fullpage}
\usepackage{color}
 
\definecolor{darkWhite}{rgb}{0.94,0.94,0.94}
 
\lstset{
    backgroundcolor=\color{darkWhite},
    breakatwhitespace=false,
    breaklines=true,
    captionpos=b,
    commentstyle=\color{red},
    deletekeywords={...},
    escapeinside={\%*}{*)},
    extendedchars=true,
    keepspaces=true,
    keywordstyle=\color{blue},
    language=python,
    morekeywords={*,...},
    numbers=left,
    numbersep=10pt,
    numberstyle=\tiny\color{black},
    showspaces=false,
    showstringspaces=false,
    showtabs=false,
    stepnumber=1,
    stringstyle=\color{gray},
    tabsize=4,
    title=\lstname,
}

%====================== INFORMATION ET RÈGLES ======================

%rajouter les numérotation pour les \paragraphe et \subparagraphe
\setcounter{secnumdepth}{4}
\setcounter{tocdepth}{2}

\hypersetup{							% Information sur le document
pdfauthor = {Premier Auteur,
			Deuxième Auteur,
			Troisième Auteur,
    		Quatrième Auteur,
    		Cinquième Auteur,
    		Sixième Auteur},			% Auteurs
pdftitle = {Nom du Projet -
			Sujet du Projet},			% Titre du document
pdfsubject = {Mémoire de Projet},		% Sujet
pdfkeywords = {Tag1, Tag2, Tag3, ...},	% Mots-clefs
pdfstartview={FitH}}					% ajuste la page à la largeur de l'écran


\usepackage{xcolor}

%======================== DÉBUT DU DOCUMENT ========================

\begin{document}

\newcommand{\E}{\mathrm{E}}
\newcommand{\Var}{\mathrm{Var}}
\newcommand{\Cov}{\mathrm{Cov}}
\newcommand{\blue}{\color{blue}}
\newcommand{\red}{\color{red}}
\newcommand{\black}{\color{black}}
\newcommand{\orange}{\color{orange}}
\newcommand{\green}{\color{green}}

%régler l'espacement entre les lignes
\newcommand{\HRule}{\rule{\linewidth}{0.5mm}}

%page de garde
\begin{titlepage}
\begin{center}

% Upper part of the page. The '~' is needed because only works if a paragraph has started.
\includegraphics[width=0.15\textwidth]{logo.jpg}~\\[1cm]

% Title
\HRule \\[0.4cm]

{
\bfseries 
\huge{C++ Pricer} \\[0.4cm]
\LARGE{Exotic products and derivatives}\\[0.4cm]}
\HRule \\[0.5cm]

% Author and supervisor
\begin{minipage}{0.4\textwidth}
\begin{flushleft} \large
\textsc{Gourio} Charles-Auguste\\

\end{flushleft}
\end{minipage}
\begin{minipage}{0.4\textwidth}
\begin{flushright} \large
\textsc charlesaugustegourio@gmail.com\\
\end{flushright}
\end{minipage}


% Bottom of the page
\vspace{0.5 in }
{\large 2024}

\end{center}
\end{titlepage}
\end{document}